% ════════════════════════════════════════════════════════════════
%  APPENDICES
% ════════════════════════════════════════════════════════════════
\appendix

\chapter{Key Code Snippets}

\section{Composite Disease Service --- Strategy Routing}

\begin{lstlisting}[language=C,caption={CompositeDiseaseService.cs --- Strategy routing and class restriction},label={lst:composite}]
public async Task<PredictionResponse> PredictAsync(PredictionRequest request)
{
    // 1. Validate image quality + content
    var validation = await _validationService.ValidateAsync(request);
    if (!validation.IsValid)
    {
        return new PredictionResponse
        {
            Label = "Rejected",
            Confidence = 0,
            IsRejected = true,
            RejectionReason = validation.RejectReason
        };
    }

    // 2. Route to correct AI service based on type
    PredictionResponse result;
    if (request.Type == DiseaseType.Pest)
        result = await _pestService.PredictAsync(request);
    else if (request.Type == DiseaseType.Weed)
        result = await _weedService.PredictAsync(request);
    else
        result = await _leafService.PredictAsync(request);

    // 3. Map to Allowed Classes
    if (isExternalApi && !result.IsRejected)
    {
        var mappedLabel = AllowedClasses.MapLabel(
            result.Label, request.Type);
        if (mappedLabel != null)
            result.Label = mappedLabel;
        else
        {
            result.IsRejected = true;
            result.Label = "Unrecognized Domain";
        }
    }
    return result;
}
\end{lstlisting}


\section{Chatbot RAG Pipeline --- Confidence Scoring}

\begin{lstlisting}[language=Python,caption={chat\_service.py --- Three-tier confidence scoring},label={lst:chat_confidence}]
HIGH_CONFIDENCE = 0.65
MEDIUM_CONFIDENCE = 0.30

def process_message(self, message, session_id=None, 
                    latitude=None, longitude=None):
    results = self.embedding_service.search(message, top_k=3)
    
    if not results:
        return self._fallback_response()
    
    top_entry, top_score = results[0]
    
    # High confidence - direct expert answer
    if top_score >= HIGH_CONFIDENCE:
        response = self._high_confidence_response(
            results, top_entry, top_score, session_id)
    # Medium confidence - partial match
    elif top_score >= MEDIUM_CONFIDENCE:
        response = self._medium_confidence_response(
            results, top_score, session_id)
    # Low confidence - out of domain
    else:
        response = self._low_confidence_response(session_id)
    
    return response
\end{lstlisting}


\section{FAISS Vector Indexing}

\begin{lstlisting}[language=Python,caption={embedding\_service.py --- FAISS index construction},label={lst:faiss}]
def _build_faiss_index(self):
    self.model = SentenceTransformer('all-MiniLM-L6-v2')
    self.embeddings = self.model.encode(
        self.search_texts, show_progress_bar=True)
    
    # Normalize for cosine similarity
    faiss.normalize_L2(self.embeddings)
    
    # Inner product on normalized vectors = cosine similarity
    dimension = self.embeddings.shape[1]
    self.index = faiss.IndexFlatIP(dimension)
    self.index.add(self.embeddings)
\end{lstlisting}


\section{Proximity Alert Generation}

\begin{lstlisting}[language=C,caption={AlertService.cs --- Proximity alert creation with Haversine distance},label={lst:alerts}]
public async Task CreateProximityAlertsAsync(
    DiseaseRecord detection, string severity)
{
    var radiusMeters = _alertSettings.RadiusInKm * 1000;
    
    var nearbyFarmers = await _userRepository
        .GetNearbyFarmersAsync(longitude, latitude, radiusMeters);
    
    // Exclude the farmer who detected the disease
    var farmersToAlert = nearbyFarmers
        .Where(f => f.Id != detection.UserId && f.Location != null)
        .ToList();
    
    var alerts = farmersToAlert.Select(farmer => new Alert
    {
        DiseaseName = detection.PredictedLabel,
        DistanceKm = CalculateDistanceKm(
            latitude, longitude,
            farmer.Location.Coordinates.Latitude,
            farmer.Location.Coordinates.Longitude),
        Severity = severity,
        IsRead = false
    }).ToList();
    
    await _context.Alerts.InsertManyAsync(alerts);
}
\end{lstlisting}


\section{Location-Aware Chatbot Response}

\begin{lstlisting}[language=Python,caption={chat\_service.py --- Location-aware contextual enrichment},label={lst:location_aware}]
# Check for nearby diseases from MongoDB
nearby_diseases = self.db_service.get_nearby_diseases(
    float(latitude), float(longitude))

# Scenario 1: User asked about a specific nearby disease
if matched_disease:
    alert = f"We have recently detected {formatted_disease} "
            f"within 5km of your location."
    response['reply'] = alert + response['reply']
    response['action'] = {
        "type": "NAVIGATE_TO_MAP",
        "params": {
            "diseaseName": matched_disease,
            "latitude": float(latitude),
            "longitude": float(longitude)
        }
    }

# Scenario 2: General location query with nearby diseases
elif is_location_query and nearby_diseases:
    disease_list = ", ".join(nearby_diseases)
    alert = f"Detected issues within 5km: {disease_list}"
    response['reply'] = alert + response['reply']
\end{lstlisting}

\newpage


\chapter{API Endpoint Documentation}

\section{Disease Detection Endpoints}

\begin{table}[H]
\centering
\caption{Disease Detection API endpoints}
\begin{tabularx}{\textwidth}{|l|l|X|}
\hline
\textbf{Method} & \textbf{Endpoint} & \textbf{Description} \\
\hline
POST & /api/disease/detect & Accepts multipart image + type (0/1/2). Returns prediction. \\
\hline
GET & /api/disease/history & Returns last 20 detections for authenticated user. \\
\hline
GET & /api/disease/map-data & Returns geolocated detections for map display. \\
\hline
GET & /api/alert & Returns all alerts for authenticated farmer. \\
\hline
GET & /api/alert/unread-count & Returns count of unread alerts. \\
\hline
PUT & /api/alert/\{id\}/read & Marks a specific alert as read. \\
\hline
\end{tabularx}
\end{table}

\section{Chatbot Endpoints}

\begin{table}[H]
\centering
\caption{Chatbot API endpoints}
\begin{tabularx}{\textwidth}{|l|l|X|}
\hline
\textbf{Method} & \textbf{Endpoint} & \textbf{Description} \\
\hline
POST & /api/chat & Send message with optional GPS. Returns AI response. \\
\hline
GET & /api/chat/welcome & Returns welcome message with categories. \\
\hline
GET & /api/chat/topics & Returns all knowledge topics grouped by category. \\
\hline
GET & /api/chat/health & Health check with knowledge base stats. \\
\hline
\end{tabularx}
\end{table}

